\documentclass[../main.tex]{subfiles}
\begin{document}

Chương 3 đã trình bày core harness AlignmentGraph-ISD như một kiến trúc năm lớp xoay quanh alignment graph. Chương 4 này mở rộng core harness đó thành một ứng dụng thực tế — case study \textbf{IT/Software Engineering ISD Prototype}: một lớp quản lý phiên cho phép cùng một chương trình đào tạo được tinh chỉnh qua nhiều phiên làm việc, cơ chế truy xuất tri thức theo domain, giao diện teacher-facing dự kiến, và giao thức đánh giá cho câu hỏi nghiên cứu thứ hai (RQ2) về hiệu quả tinh chỉnh nhiều lượt. Nội dung chương đi từ động lực thực hiện case study, qua kiến trúc ứng dụng và vòng đời phiên làm việc, qua cơ chế truy xuất tri thức miền IT/Software Engineering, đến giao diện prototype và giao thức đánh giá; kết quả đánh giá cho cả hai câu hỏi nghiên cứu được trình bày ở Chương 5.

\section{Overview and Motivation}

Đánh giá định lượng trên ISD-Agent-Bench (Chương 5) xử lý mỗi scenario độc lập trong một lượt chạy mới (\texttt{pure\_new\_run}). Cách đánh giá này phù hợp để so sánh chất lượng artifact và trajectory giữa nhiều agent trên cùng điều kiện — cơ sở để trả lời RQ1 — nhưng không phản ánh được hai đặc điểm của bài toán đặt ra ở Chương 1: một chương trình đào tạo thật được cải tiến qua nhiều phiên làm việc liên tiếp, và một miền cụ thể như IT/Software Engineering có tài liệu, quy ước và ràng buộc riêng mà một agent tổng quát khó nắm bắt đầy đủ chỉ qua system prompt. Case study, vì vậy, đảm nhận hai vai trò: cung cấp môi trường nhiều phiên để đo hiệu quả tinh chỉnh cho RQ2 bằng các process metric lấy từ chính trajectory và graph version mà harness đã ghi, và minh hoạ AlignmentGraph-ISD vận hành như một công cụ hỗ trợ thiết kế dạy học trong một ứng dụng thực tế với một domain cụ thể. Case study không nhằm thay thế benchmark hay so sánh thống kê với các agent khác.

\section{App Architecture and Session Lifecycle}

Ứng dụng case study bọc quanh core harness đã trình bày ở Chương 3 bằng một lớp quản lý phiên làm việc, cho phép một \texttt{AlignmentGraph} tồn tại lâu dài như một living artifact cấp chương trình đào tạo (training program), thay vì chỉ là working state tạm thời của một lượt chạy. Phiên đầu tiên tạo graph mới cho một chương trình đào tạo; các phiên sau đó tải lại đúng phiên bản graph đã lưu, áp dụng yêu cầu tinh chỉnh (tuning request) của người dùng, cho phép cascade invalidation lan tới các node downstream bị ảnh hưởng, rồi lưu lại thành phiên bản kế tiếp:

\begin{verbatim}
Session 1: init AlignmentGraph v1
Session 2..N: load vN, apply user tuning request,
              cascade stale downstream nodes, commit vN+1
\end{verbatim}

Để phân biệt hai chế độ trên, runtime của ứng dụng dùng hai tham số cấu hình bổ sung so với benchmark: \texttt{training\_program\_id} làm khoá bền vững của graph, và \texttt{session\_id} chỉ được truyền khi muốn tiếp tục working memory và lịch sử message của một phiên cụ thể, như minh hoạ ở Đoạn mã~\ref{lst:memory-config}.

\begin{verbatim}
# App: mo phien moi tren graph chuong trinh dao tao da co.
MemoryConfig(
    training_program_id="java-oop-training-v1",
)

# App: tiep tuc mot phien lam viec dang dang do.
MemoryConfig(
    training_program_id="java-oop-training-v1",
    session_id="tuning-session-002",
)
\end{verbatim}
\captionof{table}{Cấu hình \texttt{MemoryConfig} cho hai kiểu phiên làm việc trong ứng dụng.}
\label{lst:memory-config}

Harness tự động tải graph mới nhất nếu \texttt{training\_program\_id} đã tồn tại, tự sinh \texttt{session\_id} cho phiên mới, và tự gán loại phiên (\texttt{init} hoặc \texttt{tuning}) theo việc graph đó đã có phiên bản trước hay chưa. Khi cả hai tham số cùng được truyền, harness thêm chế độ \texttt{resume\_session} và tải lại working memory của đúng phiên đó, cho phép người dùng tiếp tục một cuộc trò chuyện tinh chỉnh khoá học đang dang dở thay vì phải mô tả lại toàn bộ ngữ cảnh từ đầu.

\section{Domain Retrieval Agents for IT/Software Engineering}

Ứng dụng case study bật cơ chế agentic RAG đã trình bày ở Phần~\ref{sec:agentic-rag} (Chương 3) và instantiate nó cho miền IT/Software Engineering: mỗi chuyên đề được đóng gói thành một retrieval subagent và đăng ký vào \texttt{RetrievalSubagentRegistry}. Cơ chế này chỉ được bật khi cấu hình \texttt{use\_rag=True} (preset \texttt{case\_study}) và có ít nhất một subagent được đăng ký; trong khi đó, benchmark ở Chương 5 luôn chạy ở chế độ \texttt{no\_rag}, nên \texttt{RetrievalPlanner} trả về kết quả rỗng và mọi agent chỉ dùng context chuẩn bị sẵn cùng graph-context/memory retrieval.

Ví dụ minh hoạ, khi ứng dụng nhận yêu cầu thiết kế một khoá Software Testing cho fresher, registry chọn subagent corpus Testing, gồm rubric test plan, đề bài unit test và ví dụ thiết kế test case, để cấp cho \texttt{AssessmentDesignAgent} và \texttt{ActivityMaterialAgent}; ngược lại, với một brief về Java/OOP, subagent Testing không được chọn mà subagent OOP được chọn thay thế. Cơ chế chọn subagent theo domain này cho phép cùng một core harness tổng quát ở Chương 3 vẫn sinh ra nội dung bám sát thực hành của từng chuyên đề IT/Software Engineering cụ thể, mà không cần huấn luyện lại hay viết prompt riêng cho từng domain.

\section{Prototype Interface}

\textcolor{red}{Trạng thái hiện tại: giao diện prototype teacher-facing dưới đây là thiết kế dự kiến (planned). Core harness (public API \texttt{MetaAgent.run(DesignBrief)}, CLI smoke test, output \texttt{addie\_output}/\texttt{trajectory}/\texttt{metadata}) đã sẵn sàng, nhưng UI teacher-facing và Markdown/PDF export chưa được hiện thực trong repo tại thời điểm viết luận văn.}

Giao diện tối thiểu dự kiến của prototype gồm sáu chức năng chính, được trình bày theo trình tự sử dụng: người dùng nhập yêu cầu thiết kế hoặc cải tiến một module đào tạo (i); chọn subject area như Java/OOP, Data Structures, Software Testing, Git, Agile/Scrum hoặc Software Engineering fundamentals (ii); chọn thời lượng, đối tượng học viên, hình thức triển khai và ràng buộc tài nguyên (iii); xem artifact ADDIE đã được sinh ra (iv); xem cảnh báo của verifier cùng lịch sử repair (v); và xuất kết quả theo hai định dạng, JSON đúng schema của benchmark để phục vụ đánh giá tự động, và Markdown/PDF lesson package cho người dạy sử dụng trực tiếp (vi). Song song với các artifact trên, giao diện còn dự kiến hiển thị workflow trace, gồm các bước agent đã thực hiện, nguồn retrieval đã dùng, kết quả validation/verifier, và chi phí token/latency của từng bước.

Prototype phục vụ hai mục tiêu bổ sung cho nhau. Thứ nhất, nó chứng minh AlignmentGraph-ISD có thể vận hành như một sản phẩm thực tế chứ không chỉ là một script chạy trên benchmark. Thứ hai, nó cung cấp một case study để phân tích sâu các artifact trong miền phần mềm, nơi yêu cầu thiết kế thường có nhiều ràng buộc domain-specific như bài tập code, rubric chấm correctness/design/readability, môi trường lab, test case và vòng lặp phản hồi.

\section{Case Study Evaluation Plan}

Case study dự kiến sử dụng ba đến năm tác vụ thiết kế hoặc cải tiến khoá học trong miền IT/Software Engineering, gồm: thiết kế module OOP Java cho fresher trong hai tuần; cải tiến bài tập Data Structures để tăng độ rõ của rubric; tạo lesson package cho Software Testing gồm lab, checklist, rubric và formative quiz; thiết kế mini-project Git/Agile cho nhóm mới vào dự án; và chuyển một workshop Code Review thành blended learning module. Mỗi tác vụ được đo bằng bốn process metric ở Bảng~\ref{tab:tuning-metrics} để trả lời RQ2, và được phân tích bổ trợ theo năm tiêu chí định tính ở Bảng~\ref{tab:case-study-criteria}.

Trọng tâm định lượng của case study là trả lời RQ2. Mỗi tác vụ được chạy thành một phiên khởi tạo (\texttt{init}) và một đến hai phiên tinh chỉnh (\texttt{tuning}) với yêu cầu chỉnh sửa cụ thể, ví dụ siết rubric, đổi thời lượng hoặc thay hình thức triển khai. Toàn bộ dữ liệu đo lấy từ những gì harness đã ghi sẵn: trajectory, \texttt{repair\_history}, các sự kiện cascade, thống kê token/latency từng bước, và diff giữa hai phiên bản graph liên tiếp; không cần bổ sung cơ chế đo mới. Bốn process metric ở Bảng~\ref{tab:tuning-metrics} thao tác hoá khái niệm ``tinh chỉnh hiệu quả'': một yêu cầu chỉnh sửa được khoanh vùng ảnh hưởng (cascade scope), chỉ sinh lại đúng phần liên quan (edit locality), với chi phí thấp hơn sinh lại toàn bộ khoá học (cost per tuning round), và không làm suy giảm alignment của các phần còn lại (alignment preservation).

\begin{table}[H]
\centering
\small
\caption{Bốn process metric đo hiệu quả tinh chỉnh nhiều lượt (RQ2).}
\label{tab:tuning-metrics}
\begin{tabular}{p{3.2cm}p{5.6cm}p{3.9cm}}
\hline
\textbf{Metric} & \textbf{Định nghĩa} & \textbf{Nguồn dữ liệu} \\
\hline
Cascade scope & Số node bị đánh dấu stale trên tổng số node sau mỗi tuning request & Cascade events trong trajectory \\
Edit locality & Tỉ lệ node được sinh lại so với node giữ nguyên giữa hai phiên bản graph & Graph version diff, \texttt{repair\_history} \\
Cost per tuning round & Token và latency của phiên tuning so với phiên khởi tạo & Thống kê token/latency từng bước \\
Alignment preservation & Alignment coverage trước và sau phiên tuning & Verifier report, coverage audit \\
\hline
\end{tabular}
\end{table}

\begin{table}[H]
\centering
\small
\caption{Tiêu chí đánh giá định tính cho case study.}
\label{tab:case-study-criteria}
\begin{tabular}{p{3.2cm}p{9.5cm}}
\hline
\textbf{Tiêu chí} & \textbf{Câu hỏi đánh giá} \\
\hline
Domain relevance & Artifact có đúng với thực hành IT/Software Engineering không? \\
Instructional alignment & Objectives, activities, materials và assessment có khớp nhau không? \\
Actionability & Giảng viên có thể dùng hoặc chỉnh sửa ít để triển khai ngay không? \\
Trace usefulness & Trace có giúp hiểu vì sao artifact được tạo hoặc sửa như vậy không? \\
Prototype usability & Workflow có phù hợp với thao tác thiết kế/cải tiến khoá học nội bộ không? \\
\hline
\end{tabular}
\end{table}

Phần đánh giá này dự kiến dùng expert review quy mô nhỏ với hai đến ba người có kinh nghiệm đào tạo nội bộ hoặc software engineering education, đối chiếu độc lập với các tiêu chí ở Bảng~\ref{tab:case-study-criteria}. Nếu chưa đủ điều kiện tổ chức human evaluation, case study sẽ được báo cáo như một phân tích định tính (qualitative analysis) dựa trên chính artifact và trace do hệ thống sinh ra, không dùng để khẳng định ưu thế thống kê so với các agent khác; vai trò của phần định tính này, vì vậy, luôn bổ sung cho các process metric của RQ2 và kết quả benchmark ở Chương 5 chứ không thay thế.

Chương này đã trình bày cách AlignmentGraph-ISD được mở rộng từ core harness sang một ứng dụng thực tế cho miền IT/Software Engineering, qua kiến trúc quản lý phiên nhiều lượt, cơ chế truy xuất tri thức theo domain, giao diện prototype dự kiến, và giao thức đánh giá gồm bốn process metric cho RQ2 cùng các tiêu chí định tính có chuyên gia. Chương 5 tiếp theo trình bày kết quả đánh giá cho cả hai câu hỏi nghiên cứu: kết quả định lượng trên ISD-Agent-Bench cho RQ1, và hiện trạng đo lường theo giao thức của chương này cho RQ2.

\end{document}
